\newcommand{\nAnalysisClasses}{17}
\newcommand{\nDesignClasses}{67}
\newcommand{\nJavaTypes}{58}
\newcommand{\nViews}{9}
\newcommand{\nTables}{17}
\newcommand{\nFks}{18}
\newcommand{\nReqs}{22}
\newcommand{\nFR}{14}
\newcommand{\nNFR}{8}
\newcommand{\nUseCases}{14}
\newcommand{\nDeps}{46}
\newcommand{\nDomainTypes}{20}


\chapter*{Abbreviations}
\addcontentsline{toc}{chapter}{Abbreviations}
\markboth{Abbreviations}{}
\begin{longtable}{@{}P{3.0cm}P{12.0cm}@{}}
\toprule
\textbf{Abbreviation} & \textbf{Meaning}\\
\midrule
\endhead
AI & Artificial Intelligence\\
API & Application Programming Interface\\
ASR & Automatic Speech Recognition (speech-to-text)\\
DB & Database\\
DDL & Data Definition Language, the SQL statements that define tables\\
ERD & Entity Relationship Diagram\\
FK / PK & Foreign key / Primary key\\
FR / NFR & Functional / Non-functional Requirement\\
HTTP / JSON & Web protocol / data format exchanged between the browser and the server\\
IS & Information System\\
JDBC & Java Database Connectivity\\
LLM & Large Language Model\\
MoSCoW & Must -- Should -- Could -- Won't, a way to rank requirements\\
ORM & Object-Relational Mapping\\
SRS & Software Requirements Specification\\
UC & Use case\\
UI & User Interface\\
UML & Unified Modeling Language\\
USD & United States dollar, the currency of all prices in the demo\\
vLLM & Open-source server that runs large language models on a GPU\\
VP & Visual Paradigm, the modelling tool used for every diagram\\
\bottomrule
\end{longtable}

\chapter*{Introduction}
\addcontentsline{toc}{chapter}{Introduction}
\markboth{Introduction}{}

\textbf{Context.} Buying a painting online is not like buying an item with a clear product code. People
remember how a painting feels rather than what it is called: ``a blue landscape'', ``Monet, but under
100 dollars'', or just a photo of a print they saw in a caf\'e. A shop that only offers a search box over
titles loses these customers. On the other hand, adding artificial intelligence without careful analysis
can make the system recommend sold-out prints, show one customer another customer's order, or stop working
whenever the AI service is slow.

\textbf{Topic.} This group essay analyses, designs and builds \textbf{Họa.}, an online gallery that sells
fine-art prints of public-domain paintings together with frames. Customers can find paintings with keywords
or a natural sentence, by speaking, with a photo, or with a photo of their own room; they can ask an AI
curator, add prints to the cart, place an order and follow it until delivery. The AI part uses the
\textbf{Gemma 4} model served by vLLM on the group's L40S server through an OpenAI-compatible interface, and
a \textbf{self-hosted Whisper} model that turns speech into text on the laptop itself. Everything else
(web pages, Java application, SQLite database) runs on one machine, so the demo works without any cloud
service. All prices are in US dollars.

\textbf{Objectives.} The essay follows the five chapters set by the lecturer, and each chapter produces a
concrete artefact of Họa.: Chapter~1 (Requirement Determination) states \nFR{} functional and \nNFR{}
non-functional requirements and \nUseCases{} use cases from the business problems of the gallery; Chapter~2
(Analysis) specifies the main use cases and models the business with activity diagrams, a domain model, a sequence
diagram and a state machine; Chapter~3 (AI in E-commerce) explains the five AI features, how each can fail and how
they are measured; Chapter~4 (Design) follows the lecturer's outline 4.1--4.6 from the analysis class diagram to the
generated code; Chapter~5 (Code) presents the source code, the user interface, the tests and how to run the demo.

\textbf{Method and consistency.} All diagrams are drawn in one editable Visual Paradigm project~\cite{visualparadigm}. The names of
requirements, use cases, classes, methods and tables are declared once in the specification file
\code{tools/model.py}; from this file the group generates the VP diagrams, the SQL schema, the Java skeleton
and the tables of this essay, and a checker compares the real code and database with it. Class and method
names follow the code in the course slides: Assignment~03 (\code{Cart}, \code{CartItem}, \code{Order},
\code{OrderItem} and the checkout transformation), Assignment~05 (Controller -- Service -- Domain --
Repository and \code{CheckoutService.checkout}) and Assignment~06 (\code{SearchUI}, \code{QueryService},
\code{SpeechService}, \code{SearchService}, \code{RankingService}, \code{ProductRepository},
\code{OrderRepository}, \code{VectorIndex}). Appendix~A maps every chapter, figure and code excerpt to its
file in the project folder.

\textbf{Scope.} Họa. is a study project. Payment and delivery are simulated; painting images come from
Wikimedia Commons~\cite{wikimedia} and only public-domain works are used; evaluation figures describe the group's own data
and test queries, not the performance of a real shop.
